\begin{exercice*}
On considère la roue de loterie suivante.
\begin{center}
\mplibforcehmode
\begin{mplibcode}
% Roue de jeu - 8 secteurs de 45°
u := 60pt;
% Couleurs RGB
color col[];
col[0] := (0.90, 0.52, 0.15); % orange
col[1] := (0.92, 0.82, 0.15); % jaune 
col[2] := (0.78, 0.58, 0.72); % rose  
col[3] := (0.90, 0.52, 0.15); % orange
col[4] := (0.18, 0.42, 0.78); % bleu  
col[5] := (0.85, 0.18, 0.18); % rouge 
col[6] := (0.78, 0.58, 0.72); % rose  
col[7] := (0.85, 0.18, 0.18); % rouge 
% Les lables 
string lbl[];
lbl[0] := "\small\bfseries 100";
lbl[1] := "\small\bfseries 10";
lbl[2] := "\small\bfseries 5";
lbl[3] := "\small\bfseries 100";
lbl[4] := "\small\bfseries 5 000";
lbl[5] := "\small\bfseries 0";
lbl[6] := "\small\bfseries 5";
lbl[7] := "\small\bfseries 0";
 
vardef drawSector(expr cx, cy, r, angA, angB, cols) =
  save p, k, da, N;
  path p;
  numeric k, da, N;
  N  := 36;
  da := (angB - angA) / N;
  p  := (cx, cy) -- (cx + r*cosd(angA), cy + r*sind(angA));
  for k = 1 upto N:
    p := p -- (cx + r*cosd(angA + k*da), cy + r*sind(angA + k*da));
  endfor;
  p := p -- cycle;
  fill p withcolor cols;
  draw p withpen pencircle scaled 1.2 withcolor white;
enddef;
% Secteurs et labels 
numeric angA, angB, amid;
pair lpos; 
for i = 0 upto 7:
  angA := 90 - (i+1)*45;
  angB := 90 - i*45;
  drawSector(0, 0, u, angA, angB, col[i]);
  amid := (angA + angB) / 2;
  lpos := (0.72*u*cosd(amid), 0.72*u*sind(amid));
  picture pic;
  pic := thelabel(lbl[i], origin);
  draw pic rotated (amid - 90) shifted lpos withcolor white;
endfor;
% Bordure 
draw fullcircle scaled (2u) withpen pencircle scaled 2;
% Cercle central 
fill fullcircle scaled 14pt;
draw fullcircle scaled 14pt withpen pencircle scaled 1.2 withcolor white;
\end{mplibcode}

\end{center}

\begin{enumerate}
    \item Si on s'intéresse aux couleurs de chaque secteur de la roue, cite les issues possibles.
    \item Si on s'intéresse aux nombres de chaque secteur, cite les issues possibles.
    \item Cite un événement certain de se réaliser.
    \item Cite un événement impossible.
    \item Quelle est la probabilité d'obtenir un nombre supérieur ou égal à 100 ?
    \item Écris à l'aide d'une phrase non négative l'événement contraire de l'événement « Obtenir un nombre supérieur ou égal à 100. »
    \item Déduis-en la probabilité de cet événement.
\end{enumerate}
\end{exercice*}
\begin{corrige}
    %\setcounter{partie}{0} % Pour s'assurer que le compteur de \partie est à zéro dans les corrigés
    %\phantom{rrr}    
    \begin{enumerate}
      \item Les issues possibles sont : orange, jaune, rose, bleu ou rouge.
      \item Les issues possibles sont : 0, 5, 10, 100 ou 5 000.
      \item Un événement certain de se réaliser est : « Obtenir une couleur. »
      \item Un événement impossible est : « Obtenir un nombre négatif. »
      \item La probabilité d'obtenir un nombre supérieur ou égal à 100 est de $\dfrac{3}{8}$, car il y a 3 secteurs avec un nombre supérieur ou égal à 100 (deux secteurs avec 100 et un secteur avec 5 000) sur les 8 secteurs au total.
      \item L'événement contraire de « Obtenir un nombre supérieur ou égal à 100 » est : « Obtenir un nombre inférieur à 100. »
      \item La probabilité de cet événement est de $\dfrac{5}{8}$, car il y a 5 secteurs avec un nombre inférieur à 100 (deux secteurs avec 0, un secteur avec 5 et deux secteurs avec 10) sur les 8 secteurs au total.
    \end{enumerate}
\end{corrige}

